% !TEX root = ../../main.tex

\section{Online Methods}

\subsection{Software implementation}
\label{sec:software-implementation}

The manuscript documents the \texttt{kl\_clustering\_analysis} software version
\texttt{0.1.0}, which is the version identifier currently exposed by the
repository. The reference implementation is written in Python and organizes the
workflow into tree construction, node-level distribution summaries,
child-parent and sibling testing, multiple-testing correction, and top-down
traversal. Version-specific engineering choices are recorded here rather than
in the main method narrative. In the software described here, the sibling basis
starts with the parent PCA directions and uses orthonormal padding if the
requested dimension exceeds the available parent rank; internal APIs can expose
child PCA projections, but the default basis constructor does not use them.

\subsection{Default analysis settings}

\Cref{tab:defaults} summarizes the default settings used in the software
described here. These settings define how the tree is built, how evidence is
combined across the tree, and how final cluster decisions are made.

\begin{table}[ht]
\centering
\begin{tabularx}{\linewidth}{@{}l>{\raggedright\arraybackslash}X@{}}
\toprule
Component & Current choice \\
\midrule
Tree distance & Hamming \\
Linkage & Average linkage \\
Edge significance level & $\alpha_{\mathrm{edge}} = 0.001$ \\
Sibling significance level & $\alpha_{\mathrm{sib}} = 0.01$ \\
Edge test & Projected Wald test with parent-specific whitening \\
Edge dimension rule & Marchenko--Pastur signal count from the parent's local spectrum \\
Sibling test & Projected Wald test with continuous edge-weight deflation \\
Sibling dimension rule & Geometric mean of positive child edge dimensions \\
Sibling basis / calibration & Parent PCA head + random padding; Satterthwaite raw test; weighted $\hat c$ localized by a Gaussian kernel in log-structural-dimension space \\
Minimum spectral dimension & 2 \\
Include internal spectral rows & On \\
Edge multiple-testing correction & Hierarchical BH correction on the tree \\
Sibling multiple-testing correction & Flat BH on focal pairs \\
Pass-through traversal & On \\
\bottomrule
\end{tabularx}
\caption{Default analysis settings used in the software documented here.}
\label{tab:defaults}
\end{table}

Two aspects of these defaults deserve special emphasis. First, the two tests do not
combine information in the same way. The child-parent edge test asks whether a
child leaves the background defined by its parent. The sibling test then asks
whether the two children are different enough from one another to justify
keeping the split. Second, once those test tables exist, the final clustering
step is simple: the method walks down the tree and keeps only the splits that
remain supported after correction.

The representation layer should also be read as part of the method. Binary data
enter directly, categorical data enter through one-hot indicator expansion, and
continuous data are discretized before the same testing procedure is applied.
